\begin{pastexam}[label={px:ncku-115-prog-05}]{成大資訊聯招 115 學年度 資料結構單選題組 B（第 5--8 題）}
（單選；每小題 3 分，答錯倒扣 1 分）

B. A university plans to connect all campus buildings by fiber links. Each building is a vertex
(ID = 0, 1, \dots, n-1), and each possible fiber link is an undirected edge \texttt{(a,b)} with
installation cost w. Assume the cost w is an integer and satisfies $0 \le \mathrm{w} \le 10^9$. The
following C program is designed to find the minimum cost of installing connections from the given links.
When two edges have the same cost, the program orders them by \texttt{(a,b)} in increasing order (smaller
\texttt{a} first; if \texttt{a} is the same, smaller \texttt{b} first). You must answer
\textbf{Questions 5--8} based on the C program below.
\begin{minted}[fontsize=\small]{c}
typedef struct { int a, b, w; } E;
static int *P, *S;
static void f0(int n){
    P = (int*)malloc(n * sizeof(int));
    S = (int*)malloc(n * sizeof(int));
    for(int i=0;i<n;i++){ P[i]=i; S[i]=1; }
}
static int f1(int x){
    while(P[x]!=x){ P[x]=P[P[x]]; x=P[x]; }
    return x;
}
static int f2(int x, int y){
    x=f1(x); y=f1(y);
    if(x==y) return 0;
    if(S[x]<S[y]){ int t=x; x=y; y=t; }
    P[y]=x; S[x]+=S[y];
    return 1;
}
static int c0(const void *px, const void *py){
    const E *x=(const E*)px, *y=(const E*)py;
    if(x->w != y->w) return x->w - y->w;
    if(x->a != y->a) return x->a - y->a;
    return x->b - y->b;
}
int main(void){
    int n, m;
    if(scanf("%d %d", &n, &m)!=2) return 0;

    E *e = (E*)malloc(m * sizeof(E));
    for(int i=0;i<m;i++){
        scanf("%d %d %d", &e[i].a, &e[i].b, &e[i].w);
        if(e[i].a > e[i].b){
            int t=e[i].a; e[i].a=e[i].b; e[i].b=t;
        }
    }
    qsort(e, m, sizeof(E), c0);
    f0(n);
    int *U = (int*)malloc((n-1) * sizeof(int));
    int *V = (int*)malloc((n-1) * sizeof(int));
    int k=0;
    long long total=0;
    for(int i=0;i<m && k<n-1;i++){
        if(f2(e[i].a, e[i].b)){
            U[k]=e[i].a;
            V[k]=e[i].b;
            total += e[i].w;
            k++;
        }
    }
    printf("%lld %d\n", total, k);
    free(e); free(U); free(V); free(P); free(S);
    return 0;
}
\end{minted}
\begin{enumerate}
\item[5.] [3\%] Which description best matches what the program does to build the backbone?
\begin{enumerate}[label=(\Alph*),nosep]
\item It starts from one vertex and repeatedly adds the cheapest edge that leaves the visited vertices.
\item It sorts all edges by cost, scans them in that order, and adds an edge only if it connects two
different groups.
\item It visits vertices level by level using a queue and records a tree of visits.
\item It repeatedly selects the closest vertex based on a distance array and updates distances.
\item It computes shortest paths between every pair of vertices, then selects edges from those paths.
\end{enumerate}

\item[6.] [3\%] What is the program output for the following input?
\begin{minted}[fontsize=\small]{text}
Input:
6 11
0 1 4
0 2 3
0 3 6
1 2 5
1 4 2
2 3 7
2 5 8
3 5 1
4 5 9
1 3 10
0 4 11
\end{minted}

(A)~\texttt{16 5} \qquad (B)~\texttt{16 4} \qquad (C)~\texttt{15 5} \qquad (D)~\texttt{15 4}
\qquad (E)~\texttt{14 4}

\item[7.] [3\%] The backbone is the graph formed by the \texttt{k} chosen edges. A vertex is an
articulation point if removing it (and its incident backbone edges) disconnects the backbone. For the
dataset in Question 6, which set of vertices are articulation points in the backbone?

(A)~$\{0, 1, 4\}$ \qquad (B)~$\{0, 2, 3\}$ \qquad (C)~$\{1, 3, 5\}$ \qquad (D)~$\{0, 1, 3\}$
\qquad (E)~No articulation point

\item[8.] [3\%] The operator wants one more feature: besides the current output, also print the number
of connected components in the input graph (using all m edges). Print three integers in one line:
\begin{center}\texttt{total \ k \ c}\end{center}
where \texttt{c} is the number of connected components. Which modification is correct?

\begin{enumerate}[label=(\Alph*)]
\item Add the following code right before \texttt{printf(...)}, and change the printf format:
\begin{minted}[fontsize=\small]{c}
int c = 0;
for(int i=0;i<n;i++){
    if(f1(i) == i) c++;
}
printf("%lld %d %d\n", total, k, c);
\end{minted}
\item Add the following code right before \texttt{printf(...)}, and change the printf format:
\begin{minted}[fontsize=\small]{c}
int c = 0;
for(int i=0;i<n;i++){
    if(P[i] == -1) c++;
}
printf("%lld %d %d\n", total, k, c);
\end{minted}
\item Add the following code right before \texttt{printf(...)}, and change the printf format:
\begin{minted}[fontsize=\small]{c}
int c = 0;
for(int i=0;i<n;i++){
    if(S[i] == 1) c++;
}
printf("%lld %d %d\n", total, k, c);
\end{minted}
\item Add the following code right before \texttt{printf(...)}, and change the printf format:
\begin{minted}[fontsize=\small]{c}
int c = 0;
for(int i=0;i<n;i++){
    c += (P[i] == i);
}
printf("%lld %d %d\n", total, k, c);
\end{minted}
\item None of the above.
\end{enumerate}
\end{enumerate}
\end{pastexam}

\begin{pastsoln}
答案：5. (B)；6. (A)；7. (D)；8. (A) 和 (D) 都對（原卷是單選，見下面的說明）。

四小題分屬：第 5、6 題 Kruskal 最小生成樹（資結書第~24.2~節），第 7 題割點（第~23.5~節），第 8 題
互斥集合（第~13~章）。 整組放在本章。

\textbf{讀程式。} \texttt{P}、\texttt{S} 是互斥集合森林（資結書第~13.2~節）：\texttt{P[x]} 是 $x$ 的
父節點，根滿足 \texttt{P[x] == x}；\texttt{S[x]} 只在 $x$ 是根時有意義，是那棵樹的大小。
\begin{itemize}
\item \texttt{f0}：每個頂點自成一組。
\item \texttt{f1}：find，往上走到根；\texttt{P[x]=P[P[x]]} 讓每個經過的節點改指向祖父，是路徑減半
（路徑壓縮的一種，資結書第~13.2.3~節）。
\item \texttt{f2}：union by size（資結書第~13.2.1~節）。 兩端已在同一組就回傳 $0$；否則把小的樹接到大的
樹的根下，回傳 $1$。
\item \texttt{c0}：先比成本，相同再比 \texttt{a}、再比 \texttt{b}；讀入時已經讓 \texttt{a < b}。
\item 主迴圈：依成本由小到大看每條邊，兩端在不同組才選，選滿 $n - 1$ 條就停。
\end{itemize}

\textbf{5. 答案 (B)。} 這正是 Kruskal 演算法（資結書虛擬碼~24.2）。 (A) 是 Prim；(C) 是 BFS；(D) 是
Dijkstra（或陣列版的 Prim）；(E) 像是 Floyd--Warshall 之類的所有點對最短路徑。 程式裡沒有佇列、
距離陣列，只有排序加互斥集合。

\textbf{6. 答案 (A)。} 排序後依序處理（$\{\cdot\}$ 是當時的分組）：
\begin{center}
\small
\begin{tabular}{cccl}
\toprule
邊 $(a, b)$ & 成本 & 動作 & 之後的分組 \\
\midrule
$(3, 5)$ & 1 & 選，$k = 1$ & $\{0\}\{1\}\{2\}\{3,5\}\{4\}$ \\
$(1, 4)$ & 2 & 選，$k = 2$ & $\{0\}\{1,4\}\{2\}\{3,5\}$ \\
$(0, 2)$ & 3 & 選，$k = 3$ & $\{0,2\}\{1,4\}\{3,5\}$ \\
$(0, 1)$ & 4 & 選，$k = 4$ & $\{0,1,2,4\}\{3,5\}$ \\
$(1, 2)$ & 5 & 跳過（同一組） & \\
$(0, 3)$ & 6 & 選，$k = 5$ & $\{0,1,2,3,4,5\}$ \\
\bottomrule
\end{tabular}
\end{center}
$k = 5 = n - 1$，迴圈停止，後面的 $(2,3,7)$、$(2,5,8)$、$(4,5,9)$、$(1,3,10)$、$(0,4,11)$ 都沒看。
$\mathit{total} = 1 + 2 + 3 + 4 + 6 = 16$，輸出 \texttt{16 5}。 (B)(D)(E) 的 $k = 4$ 不可能：
圖是連通的，生成樹一定有 $5$ 條邊。 (C) 的 $15$ 是把 $(1,2,5)$ 也算進去、漏掉 $(0,3,6)$ 的結果，但
$(1,2)$ 會形成環。

\textbf{7. 答案 (D)。} 骨幹是上面選的五條邊，它是一棵樹：
\begin{center}
\begin{tikzpicture}[x=1.5cm,y=1.3cm, v/.style={circle,draw=black!65,minimum size=6mm,inner sep=0pt,font=\small},
  w/.style={font=\footnotesize,inner sep=1.5pt}]
\node[v,fill=coreFill] (0) at (2,2) {0};
\node[v,fill=coreFill] (1) at (0.6,1) {1};
\node[v] (2) at (2,1) {2};
\node[v,fill=coreFill] (3) at (3.4,1) {3};
\node[v] (4) at (0.6,0) {4};
\node[v] (5) at (3.4,0) {5};
\draw (0) -- node[w,above left]{4} (1);
\draw (0) -- node[w,left]{3} (2);
\draw (0) -- node[w,above right]{6} (3);
\draw (1) -- node[w,left]{2} (4);
\draw (3) -- node[w,right]{1} (5);
\end{tikzpicture}
\end{center}
樹裡拿掉一個\textbf{非葉}節點，它的各個鄰居就分到不同的連通塊（樹裡兩點之間只有一條路）；拿掉葉
節點，剩下的仍是一棵樹。 所以樹的割點就是度數 $\ge 2$ 的頂點。 度數：$0$ 是 $3$，$1$ 和 $3$ 是 $2$，
$2$、$4$、$5$ 是 $1$。 割點是 $\{0, 1, 3\}$（上圖藍底）。 檢查：拿掉 $0$ 剩 $\{1,4\}$、$\{2\}$、
$\{3,5\}$；拿掉 $1$ 時 $4$ 被隔開；拿掉 $3$ 時 $5$ 被隔開。 其他選項都含有葉節點（$4$、$2$、$5$）或漏了
割點。

\textbf{8. (A)、(D) 都對。} 先確認迴圈提早停止不影響答案：
\begin{itemize}
\item 圖連通時，選滿 $n - 1$ 條邊就停，此時所有頂點已經在同一組，沒看的邊也不會再合併什麼，組數 $1$
正確。
\item 圖不連通時，選不滿 $n - 1$ 條，迴圈會把 $m$ 條邊全部看完，每條邊的兩端都被合併過，組數就是連通
元件數。
\end{itemize}
所以迴圈結束後，\textbf{根的個數}就是連通元件數。
\begin{itemize}
\item (A) 對：\texttt{f1(i) == i} 表示 $i$ 自己是根。 呼叫 \texttt{f1} 會做路徑減半，但只改非根節點的
\texttt{P}，不會改變誰是根。
\item (B) 錯：\texttt{P} 裡存的都是 $0$ 到 $n-1$ 的頂點編號，從來不會是 $-1$（那是另一種寫法「根存
負的大小」的慣例），$c$ 永遠是 $0$。
\item (C) 錯：\texttt{S} 只有根的值有意義；被接到別人下面的節點，\texttt{S} 停在接上去時的值。 第 6 題的
資料裡 $2$、$4$、$5$ 的 \texttt{S} 都還是 $1$（它們是被當成小樹接上去的單點），會數出 $3$，正確答案是 $1$。
\item (D) 對：\texttt{P[i] == i} 正是「$i$ 是根」的定義，和 (A) 數的是同一件事，只是沒有順便壓縮路徑。
\item (E) 錯：(A)、(D) 都正確。
\end{itemize}
原卷是單選題，但 (A) 和 (D) 對每一組輸入都印出相同的 $c$，兩個都正確，這一題有瑕疵；作答時選 (A) 或 (D)
都應該給分。
\end{pastsoln}
